
\subsubsection{Spurious Correlations and Worst-Group Accuracy}

Neural networks trained with ERM exploit spurious correlations --- dataset artifacts where simple co-occurring features predict class labels during training but fail to generalize across subpopulations. On Waterbirds [Wah et al., 2011; Sagawa et al., 2020] and CelebA \citep{Liu2015}, spurious correlations between background/hair color and class labels cause ERM models to achieve high average accuracy (97\% and 95\% respectively) while failing catastrophically on minority groups (worst-group accuracy 72\% and 47\%) \citep{Sagawa2020}. The theoretical basis is established: Shah et al. [2020] formally prove simplicity bias; Arpit et al. [2017] show DNNs memorize easy-to-fit patterns before hard ones, producing temporal structure exploitable for minority identification.

\subsubsection{Label-Free Debiasing Methods}

Methods requiring group annotations (GroupDRO \citep{Sagawa2020}) achieve strong worst-group performance but assume the label structure one wishes to avoid. In the annotation-free regime: JTT \citep{Liu2021} identifies minority as misclassified samples in an initial ERM run (achieving 86.7\% worst-group on Waterbirds under their experimental setup); LfF \citep{Nam2020} uses relative per-sample loss online. Both succeed because per-sample loss is an effective proxy --- minority samples incur high loss where the spurious feature conflicts with the true label. However, this magnitude-based proxy conflates hard-majority samples with spurious-minority samples [Liu et al., 2021; Kirichenko et al., 2022]. DFR \citep{Kirichenko2022} achieves near-supervised performance by retraining a linear head on a balanced held-out set, but requires a labelled validation split.

Our work extends the annotation-free regime to gradient \textit{direction} as a proxy signal, testing whether directional information distinguishes spurious-minority from hard-majority samples in a way loss magnitude cannot. We note that our contribution is signal existence (ROC-AUC measurement), not downstream task performance; we do not reproduce JTT or GroupDRO baselines in our experimental setup, and the 86.7\% worst-group figure is cited from Liu et al. [2021] and not measured here.

\subsubsection{Gradient-Based Training Dynamics Analysis}

Per-sample gradient analysis has a growing role in understanding training. Katharopoulos and Fleuret [2018] use gradient norms for importance sampling. Toneva et al. [2019] track forgetting events. Swayamdipta et al. [2020] characterize data cartography using per-epoch loss and correctness signals. Most relevant: ''Bias Leaves a Gradient Trail'' [2025] demonstrates gradient probes computed at layers can identify biased samples without labels using directional gradient signals with globally stable references --- consistent with our finding that per-batch reference is the failure point. Gradient alignment has also been studied in continual learning (GEM \citep{LopezPaz2017}) and multi-task learning \citep{Yu2020}, where global or cross-task reference directions are used --- not per-mini-batch statistics.

\textbf{Our contribution:} No prior study has directly measured within-batch gradient alignment ROC-AUC as a spurious-minority detector on Waterbirds or CelebA across training epochs. We fill this gap, characterize the failure mode, and identify the reference direction correction required.

